% Romer 6.16 --- solution fragment for hardest-solutions.tex (statement not reproduced).
% Pricing solution, output effects and the mapping to Benigno's kappa re-checked by check-11-22-6-16.py.
\problema{Romer 6.16}{Flexible- and rigid-price firms}{p.~311}{Romer \S6.5--6.7; Benigno (2015) \S4}

\passo{Environment}{All variables are in logs. Aggregate demand is $y=m-p$ (eq.~(6.41) in logs). A
firm's desired price is $p^*=p+\phi y$ (eq.~(6.58) with the constant normalized to zero), that is,
$p^*=(1-\phi)p+\phi m$ (eq.~(6.61)), with $\phi>0$ so that the flexible-price equilibrium is stable
(\S6.5). Flexible-price firms set $p^f=p^*$ after observing $m$; rigid-price firms set $p^r=Ep^*$
before. $E$ is the expectation held when rigid prices are set, so a change in $m$ is
\emph{anticipated} if it is known at that moment. A fraction $q\in(0,1)$ of firms is rigid:
$p=q\,p^r+(1-q)\,p^f$. With every price flexible, $p^*=p$ forces $y=0$: natural output is $y_n=0$
and $p=m$.}

\parte{a}
\passo{Substitute the price index}{Inserting $p=q\,p^r+(1-q)\,p^f$ into $p^f=(1-\phi)p+\phi m$,
\[ p^f=(1-\phi)\big[q\,p^r+(1-q)\,p^f\big]+\phi m . \]}

\passo{Collect $p^f$}{Moving the $p^f$ terms to the left,
\[ p^f\big[1-(1-\phi)(1-q)\big]=(1-\phi)\,q\,p^r+\phi m ,\qquad
1-(1-\phi)(1-q)=q+\phi(1-q)\equiv D>0 . \]
Dividing by $D$ and using $D-\phi=(1-\phi)q$,
\[ p^f=\frac{(1-\phi)\,q\,p^r+\phi m}{q+\phi(1-q)}=p^r+\frac{\phi}{q+\phi(1-q)}\,\big(m-p^r\big). \tag{A} \]
Flexible-price firms close the fraction $\phi/D$ of the gap between money and the rigid price. Since
$D-\phi=(1-\phi)q$, $\phi/D<1$ exactly when $\phi<1$, i.e.\ when desired prices rise less than
one-for-one with the price level ($\partial p^*/\partial p=1-\phi$: strategic complementarity).}

\parte{b}
\passo{Expected flexible price}{The flexible firms' rule holds whatever $m$ turns out to be, and
expectations are linear, so $Ep^f=(1-\phi)\,Ep+\phi\,Em$. The right side is exactly the rigid
firms' rule: $Ep^f=p^r$. Each rigid firm expects to charge what the flexible firms will charge.}

\passo{Expected price level}{$p^r$ is chosen with the information behind $E$, so $Ep^r=p^r$ and
\[ Ep=q\,p^r+(1-q)\,Ep^f=q\,p^r+(1-q)\,p^r=p^r . \]}

\passo{Solve}{Substituting $Ep=p^r$ into the rigid firms' rule and dividing by $\phi>0$,
\[ p^r=(1-\phi)\,p^r+\phi\,Em\quad\Longrightarrow\quad \phi\,p^r=\phi\,Em
\quad\Longrightarrow\quad p^r=Em . \tag{B} \]
No parameter appears: rigid firms preset the price level that flexible prices would deliver given
what they know, and $Ey=Em-Ep=0$. Check against (A): with $p^r=Em$,
$Ep^f=Em+\frac{\phi}{D}(Em-Em)=Em=p^r$, as the first step requires.}

\parte{c}
\passo{Price level and output}{With (B), (A) reads $p^f=Em+\frac{\phi}{D}(m-Em)$, so
\[ \begin{aligned}
p&=q\,Em+(1-q)\Big[Em+\frac{\phi}{D}(m-Em)\Big]=Em+\frac{(1-q)\,\phi}{D}\,(m-Em),\\
y&=m-p=\Big[1-\frac{(1-q)\,\phi}{D}\Big](m-Em)=\frac{q}{q+\phi(1-q)}\,(m-Em),
\end{aligned} \tag{C} \]
because $D-(1-q)\phi=q$.}

\passo{(i) Anticipated changes}{Write $m=Em+u$, where $u\equiv m-Em$ is the part of money unknown
when rigid prices were set; by (C), $y=(q/D)\,u$ depends on $u$ alone. An anticipated change raises
$Em$ and $m$ by the same $\Delta$ and leaves $u$ unchanged: $\Delta y=0$, while $p^r$, $p^f$ and $p$
all rise by $\Delta$. Money is neutral because the rigid firms knew the change when they set their
prices, so every price incorporates it, and the flexible firms keep their relative price
$p^f-p^r=(\phi/D)\,u$ unchanged.}

\passo{(ii) Unanticipated changes}{Raising $m$ by $\Delta$ with $Em$ fixed raises $u$ by $\Delta$:
\[ \Delta y=\frac{q}{q+\phi(1-q)}\,\Delta>0,\qquad \Delta p=\frac{(1-q)\,\phi}{q+\phi(1-q)}\,\Delta<\Delta . \]
Since $D-q=\phi(1-q)>0$, $0<q/D<1$. Output absorbs part of the surprise because the fraction $q$ of
prices is stuck at $Em$; by (A), flexible firms raise their prices by $(\phi/D)\Delta$, which is less
than $\Delta$ when $\phi<1$.}

\passo{Mapping to Benigno's AS curve}{The AS curve follows from the pricing side alone. From the
price index and $Ep=p^r$, $p-Ep=(1-q)(p^f-Ep)$; and $p^f-Ep=(p^f-p)+(p-Ep)=\phi y+(p-Ep)$. Hence
\[ q\,(p-Ep)=(1-q)\,\phi\,y\quad\Longrightarrow\quad p-Ep=\frac{(1-q)\,\phi}{q}\,(y-y_n),
\qquad y_n=0 . \tag{D} \]
This is Benigno's eq.~(17), $p-p^e=\kappa\,(y-y_n)$, under the dictionary: the rigid share $q$ is
his share $\alpha$ of prices preset at $p^e$; $p^r=Ep$ is his $p^e$; $p^f$ is his $\tilde p$, whose
rule $\tilde p-p=(\sigma^{-1}+\eta)(y-y_n)$ (his eq.~(16)) is $p^*-p=\phi y$ with
$\phi=\sigma^{-1}+\eta$. Therefore $\kappa=(1-q)\phi/q=(1-\alpha)(\sigma^{-1}+\eta)/\alpha$. Romer's
\S6.5 economy is the special case $U=C-L^\gamma/\gamma$: utility is linear in $C$ ($\sigma^{-1}=0$)
and $L^\gamma/\gamma=L^{1+\eta}/(1+\eta)$ with $\eta=\gamma-1$, so $\phi=\gamma-1=\sigma^{-1}+\eta$.}

\passo{Output effect and real rigidity}{Combining (D) with the AD curve and $Ep=Em$,
$y=(m-Em)-(p-Ep)=(m-Em)-\kappa y$, so
\[ \begin{gathered}
y=\frac{m-Em}{1+\kappa},\qquad \frac{1}{1+\kappa}=\frac{q}{q+(1-q)\,\phi},\\
\frac{\partial}{\partial\phi}\Big(\frac{q}{q+\phi(1-q)}\Big)=-\frac{q\,(1-q)}{\big[q+\phi(1-q)\big]^2}<0 ,
\end{gathered} \]
which reproduces (C). Lower $\phi$ (more real rigidity) means a flatter AS (lower $\kappa$) and a
larger output response to the same surprise. Limits: as $\phi\to0$, $y\to m-Em$, so the whole
surprise goes to output although a fraction $1-q$ of firms is free to reprice (with $p^*=p$ they
choose not to); at $\phi=1$, $p^f=m$ and $y=q\,(m-Em)$; as $\phi\to\infty$, $y\to0$. With Benigno's
calibration ($\alpha=0.66$, $\sigma=0.5$, $\eta=0.2$; his \S8), $\kappa=0.34\times2.2/0.66\approx1.13$
and a 1\% money surprise raises $y$ by $1/(1+\kappa)\approx0.47\%$; keeping $\alpha=0.66$ but setting
$\phi=0.2$ gives $\kappa\approx0.10$ and $\approx0.91\%$.}

\intuicao{Real rigidity alone is neutral: with $q=0$, $y=0$ for every $\phi$ ($\kappa\to\infty$). It
magnifies nominal rigidity: when $\phi$ is small, flexible firms care mostly about their price
relative to the sticky ones, so they track the preset prices rather than money, and a small share of
preset prices freezes much of the price level (\S6.7). Anticipated money is neutral for any $q$ and
$\phi$, because no price is stuck at a stale value.}

\armadilha{(1) ``Anticipated'' refers to the information held when rigid prices are set: a change
announced after that date is a surprise here even if it is public before $m$ is realized. (2) $p^r$
is a choice, not an outcome to forecast: $Ep^r=p^r$, which is what delivers $Ep=p^r$. (3) Romer's
$\eta$ (the elasticity of demand, Benigno's $\theta$) is not Benigno's $\eta$ (the inverse Frisch
elasticity); and Benigno's $\sigma\equiv\tilde\sigma s_c$ is the intertemporal elasticity scaled by
the consumption share.}

\refbib{Romer (2012), \S6.5, eqs.~(6.41), (6.58), (6.61), and \S6.7 (real rigidity); Benigno (2015),
\S4.1, eqs.~(16)--(17), calibration in \S8.}

\checado{\texttt{check-11-22-6-16.py} solves the pricing system with sympy by brute force (without the
shortcut $Ep^f=p^r$) and confirms (A), $p^r=Em$, (C), $\partial y/\partial Em=0$ at given $u$,
$0<q/D<1$, (D) with $\kappa=(1-q)\phi/q$, its identity with $(1-\alpha)(\sigma^{-1}+\eta)/\alpha$ under
the mapping, $y=(m-Em)/(1+\kappa)$, the sign of $\partial y/\partial\phi$ and the calibration numbers.
The solution manual agrees: same $p^f$, $p^r=Em$, $y=q\,(m-Em)/[\phi+(1-\phi)q]$ and the same
real-rigidity derivative (its eq.~(19)); it does not discuss Benigno's AS.}
